% ============================================================
% Section: Future work
% Standalone section file, like sections/custom_dataset.tex.
% Requires: \usepackage{booktabs} in the main document preamble.
% ============================================================

\section{Future Work}
\label{sec:future-work}

\subsection{A Keyless Fallback Search Route}
\label{subsec:future-ddg-fallback}

\paragraph{The problem.} Every piece of web evidence the system
retrieves passes through a single self-hosted SearXNG instance, which
in turn scrapes public search engines from one IP address without API
keys. That design keeps the system free of paid services and of
credentials, but it makes the search a single point of failure. On
25~September~2026 it failed almost completely: 68 of 69 queries sent
during a measurement came back empty, because the upstream engines had
rate-limited or challenged the instance, and 23 of 24 claims were left
without any web evidence. Each of those claims was reported as
\texttt{UNVERIFIED}, indistinguishable from a claim that had been
searched and found unsupported.

Since 30~September~2026 the system reports such a failure explicitly:
when every query planned for a claim fails, the claim is marked
\texttt{searchUnavailable} and its trace stops at the evidence
retrieval stage, so the evaluation can report those claims apart from
genuine \texttt{UNVERIFIED} verdicts. This makes the failure visible;
it does not recover the evidence that was lost.

\paragraph{Why not an API-keyed engine.} The usual remedy is a search
engine accessed with an API key (for example, the Brave Search API or
Google Programmable Search), which is rate-limited per key rather than
per scraping IP. It is deliberately not adopted here: it introduces a
per-query cost, a credential that every deployment must manage, and a
dependence on a commercial service whose terms and free quotas may
change, which would also weaken the reproducibility of the results.

\paragraph{Proposed design.} A second, independent route that is
used only when the first one fails:
\begin{itemize}
    \item \textbf{Trigger on failure only.} When SearXNG cannot answer
    a query (the condition that already raises the system's
    \texttt{SearchUnavailableError}), the same query is sent once to
    DuckDuckGo directly, bypassing SearXNG and its engine-suspension
    state. In normal operation the fallback sends nothing, so it adds
    no load and no latency.
    \item \textbf{Same shape of result.} DuckDuckGo results are mapped
    to the same evidence records (URL, title, snippet) and tagged with
    their route in the existing \texttt{engines} field, so every later
    stage --- the pertinence gate, ranking and verification ---
    treats them identically, and the evaluation can still tell which
    route produced each source.
    \item \textbf{Its own ceiling.} Requests to the fallback are
    bounded by a process-wide concurrency limit and back off after a
    refusal, like every other external service the system calls; a
    per-request override of that limit is not allowed, since the load
    falls on a service shared by all concurrent runs.
    \item \textbf{A stricter definition of ``unavailable''.} A claim
    is marked \texttt{searchUnavailable} only when both routes fail.
\end{itemize}

\paragraph{The risk that must be measured first.} DuckDuckGo does not
offer a free API for web results --- its Instant Answer API returns
summary information, not a ranked list of pages --- so a keyless
route means requesting its HTML results page and parsing it. That has
two consequences. First, the parser depends on markup that can change
without notice. Second, and more seriously, the request is subject to
the same bot detection the fallback is meant to escape: in the
25~September measurement, DuckDuckGo queried through SearXNG answered
with a CAPTCHA or a timeout on every request (104 CAPTCHAs and 231
timeouts in the preceding three hours), from the same machine that
would send the fallback's requests. A fallback that fails under the
same conditions as the primary route adds nothing. The fallback is
therefore proposed together with its evaluation, and should only be
adopted if it succeeds:

\begin{table}[ht]
\centering
\caption{Measurements required before adopting the fallback route.}
\label{tab:future-ddg-measurements}
\begin{tabular}{@{}ll@{}}
\toprule
\textbf{Measurement} & \textbf{Question it answers} \\
\midrule
Answer rate under benchmark load & Does it answer when SearXNG does not? \\
Answer rate from the deployment IP & Is the result specific to one machine? \\
Share of \texttt{searchUnavailable} claims & How much evidence does it recover? \\
Accuracy on the recovered claims & Is the recovered evidence useful? \\
Added latency per recovered claim & What does the recovery cost? \\
\bottomrule
\end{tabular}
\end{table}

The first two rows decide whether the route is worth building at all;
the last three come from the evaluation harness, run once with and
once without the fallback over the same labelled claims.
